\documentclass[12pt,a4paper]{article}

\usepackage[utf8]{inputenc}
\usepackage[T1]{fontenc}
\usepackage[slovak]{babel}
\usepackage{geometry}
\usepackage{hyperref}
\usepackage{tabularx}
\usepackage{enumitem}

\geometry{margin=2.5cm}
\setlist[itemize]{topsep=2pt,itemsep=2pt}

\title{02\_login\_final -- Záverečný report}
\author{Tím Kapuut}
\date{\today}

\begin{document}

\maketitle

\section{Cieľ a zadanie projektu}
\begin{itemize}
    \item Stručné zhrnutie pôvodného zadania (funkcionality, cieľové prostredie, kľúčové obmedzenia).
    \item Zdôraznenie odchýlok oproti pôvodnému zadaniu: \textit{čo sa menilo, prečo a ako sa to prejavilo v implementácii}.
    \item Hlavné pridané hodnoty oproti zadaniu: \textit{bonusové funkcionality, zlepšenia použiteľnosti alebo výkonu}.
\end{itemize}

\section{Rozdelenie práce členov tímu na výslednej FE implementácii}
\begin{tabularx}{\textwidth}{|>{\bfseries}l|>{\bfseries}l|X|}
    \hline
    Člen tímu & Oblasť FE & Stručný popis príspevku \\
    \hline
    \textit{Meno 1} & \textit{napr. UI navigácia} & \textit{hlavné komponenty, obrazovky, refaktoring, štýly} \\
    \hline
    \textit{Meno 2} & \textit{napr. minihry} & \textit{implementované scény, správanie, integrácie} \\
    \hline
    \textit{Meno 3} & \textit{napr. stav aplikácie} & \textit{manažment dát, prepojenie so serverom, optimalizácie} \\
    \hline
    \textit{Meno 4} & \textit{napr. vizuál} & \textit{dizajnové úpravy, animácie, responzivita} \\
    \hline
\end{tabularx}

\section{Zapracovanie pripomienok z kontrolnej prezentácie}
\begin{itemize}
    \item \textbf{Pripomienka 1:} \textit{stručný opis pripomienky}.
    \item \textbf{Úprava:} \textit{čo sme zmenili, kto to riešil, aký je stav (dokončené/potrebné dopracovať)}.
    \item \textbf{Dopad:} \textit{vplyv na UX/FW, prípadné riziká}.
    \item \textbf{Pripomienka 2:} \textit{... (pokračujte pre každú pripomienku)}.
\end{itemize}

\section{Testovanie}
\subsection{Prístup a rozsah testovania}
\begin{itemize}
    \item Typy testovania: \textit{manuálne scenáre, exploratory, prípadne automatizované skripty}.
    \item Pokryté časti FE: \textit{navigácia, kľúčové scény, kritické minihry, chybové stavy}.
    \item Použité nástroje a prostredia: \textit{verzia engine, OS, zariadenia}.
\end{itemize}


\subsection{Zistené problémy a ich riešenie}
\begin{itemize}
    \item \textbf{Najdôležitejšie chyby:} \textit{popis, priorita, stav opravy}.
    \item \textbf{Regresné riziká:} \textit{oblasti s nedostatočným pokrytím alebo závislosťami}.
    \item \textbf{Plán do odovzdania:} \textit{čo ešte otestovať alebo opraviť}.
\end{itemize}


\end{document}
